\section{The estimates that remain open}
The all-simplicity question is the unboundedness of the exact finite
heights in Theorem~\ref{thm:fd-arithmetic-complete}:
\[
 \sup_{N,k}\left(\frac2{S_k-\operatorname{tr}C_N^{2k}}\right)^{1/(2k)}
          =+\infty.
\]
The arithmetic computations establish this height bound through
$97.632$. The equivalent assertion at every height remains open.
The dimension and fixed-space limits in
Proposition~\ref{prop:fd-certificate-size} identify two constraints
on a quantitative proof: its dimension must reach twice the counted
packet size, and increasing the moment order in one fixed space
eventually returns the threshold to the first ordinate. The unresolved
estimate concerns the finite arithmetic defect with both indices
allowed to grow.

Equivalently, the nonnegative integer
\(n(T)-\operatorname{ind}_+(T H_N-G_N)\) must vanish for finite
indices at unbounded zero-free heights. The logarithmic-support
construction proves a lower inertia density \(d_{\rm MT}-\eta\)
at every sufficiently large height; the complementary integer deficit
is still uncontrolled. Equation~\eqref{eq:fd-trace-inertia-integral}
transports this assertion exactly to the trace family. The rank-four
comparison in \eqref{eq:fd-inertia-surgery} preserves the density
conclusion and changes the limiting integer deficit by one.

The prime form also has the exact multiplication identity
\eqref{eq:fd-prime-commutator}, whose spectral side consists of
derivative evaluations at the zeros. Proposition~\ref{prop:fd-multiplication-domain}
constructs compact mean-zero tests tending to zero in the Weil norm
whose products by the coordinate converge to a nonzero spectral
coordinate. Thus this test operator has no closed extension, and a
cyclic trace argument through that proposed extension cannot provide
the missing divisor trace. The explicit prime form identity remains
available; a quantitative comparison retaining its derivative terms is
still unproved.

The derivative-moment formulation has an equally precise absolute
threshold. On every bounded packet, the minimum of the polynomial
square sum with constant term one is the total length of the
multiple-zero fibers. Its value below $2$ certifies simplicity.
The contour formula and coefficientwise error budget in
Proposition~\ref{prop:fd-moment} specify a finite test. Leading
asymptotic derivative moments allow errors proportional to the
number of zeros, while this test asks for a certified absolute
margin below $2$. The exact positive-kernel replacement in
Section~\ref{sec:fd-positive-surgery} preserves the stated leading
moments, simple-zero proportions, and zero heat threshold, while
its divisor and prime-form corrections are explicit. Thus the
missing step in this route must use arithmetic information beyond
those transported asymptotic quantities. The boundary identities
\eqref{eq:poisson-infimum} and \eqref{eq:prime-ramp-coefficient}
express that information as two further exact scalar thresholds;
their required upper estimates remain unproved.

For the original prolate trial, the unresolved ratio is
$\liminf q(v_L)/a_2(L)=0$. The trial upper estimate has scale
$\lambda^{9+\varepsilon}e^{-4\pi\lambda^2}$, where
$\lambda=e^L$. The finite matrix bounds of
Section~\ref{sec:finite-weil-lower} give a concrete lower certificate
with a frequency tail and a quadratic coupling error. No arithmetic
estimate establishing a sufficient positive margin on an unbounded
sequence of windows has been proved. The comparison with the odd
trial places the remaining positive polynomial gain in
$0<\alpha\le4$, meaning the absolute scale
$\lambda^{9+\alpha}e^{-4\pi\lambda^2}$.
The all-window lower rate in Theorem~\ref{thm:cosine-lower-rate}
provides \(\exp\{-C\lambda^{2/d_{\rm MT}}\log\lambda
(1+\log\log\lambda)\}\) for all three variational quantities.
Its explicit exponent comparison with the trial leaves the required
ratio estimate open.
The compact-operator example proves that the exact commutator and
the indicated compactness estimates permit a divergent quotient
norm; every transfer from them to the desired denominator must
therefore include an additional quantitative estimate for this form.

M\"obius inversion gives the exact weighted physical norm in
\eqref{eq:mobius-physical-plancherel} and the local pole masses in
\eqref{eq:mobius-abel}. Its compact smoothing of normalized Mertens
is \eqref{eq:mobius-mertens-smoothing}. The fixed positive shifts
are controlled; an \(O(\epsilon^{-1})\) squared-norm estimate as
\(\epsilon\downarrow0\) is not established. The explicit odd test
in that section also shows that the one-mode orthogonality constraint
leaves these pole evaluations present. These formulas preserve the
higher-pole and derivative factors that a sharper arithmetic estimate
must control.

For Burnol's local jets, Proposition~\ref{prop:burnol-coupled-cutoff}
computes the discarded Weil norm for a fixed polynomial and a fixed
ordinate. Its compact Rayleigh quotient is asymptotic to
$m_\gamma(A_{\chi,M}^2/B_{\chi,M})L$. The estimate uniform in
height and in the number of jet orders, needed to apply this construction
to growing packets, remains open. The boundary-kernel identity
\eqref{eq:hb-shift-scale} gives another exact expression of the
same multiplicity: at each fixed ordinate,
$\lim_{\omega\downarrow0}K_\omega(\gamma,\gamma)/\omega$
is positive precisely when $m_\gamma=1$. A positive arithmetic
lower bound for this limit at every ordinate has not been proved.

Half balance of the two arithmetic Hurwitz directions also remains
open. Proposition~\ref{prop:periodic-gaussian-gate} gives the finite
Abel estimate that isolates its missing data: fixed odd phase Fourier
means and the mass near the integer phase boundaries. The explicit
phase-measure map proved there preserves the established even Fourier
and normalized Gaussian information while fixing one sign. The required
odd-mode and boundary estimates for the arithmetic labels remain
unproved.

The contour and probe tests have now been executed on the same
positive-ordinate packet through height $20$. The contour calculation
also certifies the eight-zero packet through height $45$.
For larger contours the leading-coefficient bound needs separation
of distinct centers, while its Taylor estimate needs a zero-free
complex-time boundary and a logarithmic-derivative bound there.
The probe needs a spectral isolation bound and a residual small
enough that its divided mass has error below $1/2$.
All these quantities enter explicit inequalities in
Sections~\ref{sec:local} and~\ref{sec:suzuki-probes}.
Suzuki's convergence gives eventual recovery at every fixed
ordinate, and hence eventual simultaneous recovery on each fixed
finite packet. It supplies neither the effective index used in the
present computation nor the conclusion that every recovered integer
equals one. The all-height task is to certify that last conclusion
for packets covering every positive ordinate.

For zeta collisions, the isolated-pair expansion is proved, including
the arithmetic coordinate conversion. The statistical questions
still require control of the exterior interaction
$\delta^2\sum_{\xi\ne x_-,x_+}(b-\xi)^{-2}$ on shrinking gaps,
of nearby triples, and of the remaining simple branches. The
fixed-scale gap results and the restricted correlation identities
have their exact comparisons in
Section~\ref{sec:arithmetic-collisions}; the displayed sine-product
and separated-sequence examples prove the limitations of those
particular statistics. A zeta Weibull or Poisson collision law is
not established by the present results.
